The economic attraction of contingency is easiest to see from the timing of a public budget. A sovereign's debt service is fixed in advance, while revenues and emergency spending respond to realized economic conditions. A disaster can therefore produce a financing need on both sides of the budget: output and the tax base fall just as recovery spending becomes urgent. A contract that changes the payment due in that state can reduce the need for an immediate restructuring or abrupt fiscal adjustment. Whether this timing improvement also increases the amount raised at issuance is a separate question, because lenders anticipate the reduction in service and incorporate it into the price. The paper studies this second margin explicitly rather than assuming that every reduction in bad-state payments creates additional borrowing capacity.

This perspective connects debt management to the distinction between chronic and acute physical risk. Chronic damages change the distribution of future resources as warming accumulates. Acute shocks make some realizations of that distribution particularly adverse, even at the same conditional trend. A liability portfolio designed around average growth alone can therefore miss the states that determine market access. Adding disaster jumps preserves the tractability of a reduced-form growth model while making the repayment schedule's exposure to rare adverse states visible. It also makes clear why a local disaster indicator and an output indicator need not convey the same information about the resources available to creditors.

The practical comparison involves information and implementation as well as insurance. An output trigger observes a broad measure of economic performance, including shocks unrelated to the physical event. A disaster trigger can instead rely on a verified event or a loss assessment tied to a contractual threshold. The latter can be simpler to specify for an individual hazard, but the resulting protection depends on which losses the index captures. Output measurement, data revisions, the timing of certification, and contractual enforceability could all change the relative merits of the two designs. They are not priced in the numerical exercise, so its output-linked benchmark should be interpreted as a transparent information benchmark against which trigger mismatch can be measured.

The distinction between borrowing capacity and welfare is equally important. More valuable claims can finance a larger issue, but a sovereign may use the proceeds to increase exposure rather than to preserve a buffer. Conversely, a contract that lowers issuance proceeds can still smooth expenditure in a valuable way. Establishing either welfare ranking requires a government objective, adjustment costs, and the distribution of consumption or public services. The fiscal-limit approach isolates an earlier question: how does the repayment mapping change the maximum market value of liabilities consistent with the assumed resources and default rule? Answering it is useful for understanding contractual design without treating a proceeds gain as a complete policy evaluation.
